\section{Numerical Verification}
\label{sec:verification}

Before evaluating machine-learning-based refinement strategies, a uniform-mesh
refinement study was performed to establish the numerical behavior of the
baseline CFD model. Grid-refinement studies and observed-order analyses are
standard components of CFD solution verification
\cite{roache1998,celik2008}.

Four structured meshes were considered, denoted M20, M40, M80, and M160.
Successive refinement approximately halves the characteristic cell spacing
in each coordinate direction.

\begin{table}[htbp]
\centering
\caption{Uniform meshes used for the numerical verification study.}
\label{tab:uniform_meshes}
\begin{tabular}{lrrr}
\toprule
Mesh & Cells & $\Delta x$ [m] & Minimum $U_x/U_{\mathrm{lid}}$ \\
\midrule
M20  & 400    & 0.00500 & -0.198089 \\
M40  & 1600   & 0.00250 & -0.209591 \\
M80  & 6400   & 0.00125 & -0.212887 \\
M160 & 25600  & 0.000625 & -0.213733 \\
\bottomrule
\end{tabular}
\end{table}

A vertical velocity sampling line located at the horizontal center of the
cavity,

\begin{equation}
x/L = 0.5,
\end{equation}

was used to compare the horizontal velocity component between successive
meshes.

The resulting inter-grid differences are summarized in
Table~\ref{tab:grid_convergence}.

\begin{table}[htbp]
\centering
\caption{Successive-mesh differences for the vertical centerline velocity
profile.}
\label{tab:grid_convergence}
\begin{tabular}{lrr}
\toprule
Mesh comparison & $L_2$ difference & $L_\infty$ difference \\
\midrule
M20--M40   & $6.1826\times10^{-3}$ & $1.2598\times10^{-2}$ \\
M40--M80   & $1.7211\times10^{-3}$ & $3.5090\times10^{-3}$ \\
M80--M160  & $4.2841\times10^{-4}$ & $8.8600\times10^{-4}$ \\
\bottomrule
\end{tabular}
\end{table}

The reduction ratios of the successive \(L_2\) differences are approximately

\begin{equation}
\frac{E_{20,40}}{E_{40,80}} = 3.59,
\qquad
\frac{E_{40,80}}{E_{80,160}} = 4.02.
\end{equation}

For a refinement ratio of approximately two, an observed convergence order
can be estimated from the successive solution differences, consistent with
standard grid-refinement verification practice
\cite{roache1998,celik2008}:

\begin{equation}
p =
\frac{\ln(E_h/E_{h/2})}{\ln(2)}.
\end{equation}

The corresponding observed orders are approximately

\begin{equation}
p = 1.845
\end{equation}

for the M20--M40--M80 sequence and

\begin{equation}
p = 2.006
\end{equation}

for the M40--M80--M160 sequence.

These results indicate strong convergence of the sampled centerline velocity
profile, with behavior approaching second order on the finer meshes.

\subsection{Reference solution used in the adaptive-mesh study}

The M160 solution is subsequently used as a fine-grid numerical reference
for comparisons with the adaptive meshes. It is important to emphasize that
M160 is not treated as an exact analytical solution or as a formally
grid-independent solution.

Consequently, quantities reported later as ``error'' relative to M160 should
more precisely be interpreted as discrepancies with respect to the selected
fine-grid numerical reference.

\subsection{Verification limitations}

The time step was reduced together with the spatial mesh size during the
uniform refinement sequence. Therefore, the present refinement study does
not rigorously separate spatial discretization error from temporal
discretization effects.

Furthermore, the discontinuity between the moving lid and the stationary
side walls produces strong gradients near the upper corners of the cavity.
These regions can dominate pointwise error measures and must be considered
carefully when interpreting maximum-error metrics.

For these reasons, the verification study is used to establish numerical
consistency and to define a controlled internal reference for the subsequent
comparisons, rather than to claim an exact discretization-error estimate.
